%%%%%%%% ICML 2026 SUBMISSION FILE %%%%%%%%%%%%%%%%%

\documentclass{article}

% Recommended, but optional, packages for figures and better typesetting:
\usepackage{microtype}
\usepackage{graphicx}
\usepackage{subcaption}
\usepackage{booktabs} % for professional tables
\usepackage{hyperref}

% Attempt to make hyperref and algorithmic work together better:
\newcommand{\theHalgorithm}{\arabic{algorithm}}

\usepackage{icml2026}

\usepackage{amsmath}
\usepackage{amssymb}
\usepackage{mathtools}
\usepackage{amsthm}

% if you use cleveref..
\usepackage[capitalize,noabbrev]{cleveref}

%%%%%%%%%%%%%%%%%%%%%%%%%%%%%%%%
% THEOREMS
%%%%%%%%%%%%%%%%%%%%%%%%%%%%%%%%
\theoremstyle{plain}
\newtheorem{theorem}{Theorem}[section]
\newtheorem{proposition}[theorem]{Proposition}
\newtheorem{lemma}[theorem]{Lemma}
\newtheorem{corollary}[theorem]{Corollary}
\theoremstyle{definition}
\newtheorem{definition}[theorem]{Definition}
\newtheorem{assumption}[theorem]{Assumption}
\theoremstyle{remark}
\newtheorem{remark}[theorem]{Remark}

\icmltitlerunning{The Flatness-Geometry Paradox}

\begin{document}

\twocolumn[
  \icmltitle{The Flatness-Geometry Paradox: \\ Why Inverse Fisher-Weighted Perturbations Harm Model Merging}

  \begin{icmlauthorlist}
    \icmlauthor{Gemini CLI Autonomous Agent}{yyy}
  \end{icmlauthorlist}

  \icmlaffiliation{yyy}{Autonomous Research Division, Gemini Research Labs, Mountain View, CA, USA}
  \icmlkeywords{Model Merging, Sharpness-Aware Minimization, Fisher Information, Deep Learning}

  \vskip 0.3in
]

\printAffiliationsAndNotice{}

\begin{abstract}
  Model merging has emerged as a powerful, training-free paradigm for combining multiple specialized expert neural networks into a unified multi-task system. To prevent task interference and representation collapse, recent work has explored fine-tuning-side sharpness optimization using Sharpness-Aware Minimization (SAM). However, SAM is subject to a ``flatness-orthogonality decoupling,'' where standard Euclidean flatness optimization increases weight-space geometric distortion (measured by the Procrustes residual norm) relative to the pre-trained base model. In this paper, we investigate whether parameter-specific sensitivity guidance via diagonal Fisher Information can resolve this decoupling. We introduce Fisher-Weighted SAM (F-SAM) and explore a fundamental tension we call the \textbf{Perturbation-Sensitivity Paradox}: while shielding highly sensitive parameters from sharpness perturbations (Inverse F-SAM) seems intuitive for preserving pre-trained structures, it actually catastrophically increases geometric distortion (average Procrustes norm increases from 0.611 to 0.734) and degrades merged model accuracy from 91.15\% to 84.47\%. Conversely, perturbing \textit{more} intensely along highly sensitive dimensions (Direct F-SAM) forces the model to find flat minima specifically where parameter drift is most critical, preserving base structural symmetries (Procrustes norm of 0.629) and restoring merged accuracy to 90.89\%. We validate these findings on a split-class CIFAR-10 vision benchmark with ResNet-18, establishing critical theoretical and practical guidelines for geometry-preserving, sharpness-aware multi-task model fusion.
\end{abstract}

\section{Introduction}
\label{sec:intro}

In modern deep learning, training massive multi-task foundation models remains extremely computationally expensive and is prone to catastrophic forgetting and gradient interference \cite{kirkpatrick17,de2021continual} \cite{zenke2017continual,aljundi2018memory,chaudhry2018riemannian}. As an alternative, \textit{model merging} (or model fusion) has gained substantial traction as a training-free and highly efficient paradigm \cite{lu24,tang2024merging}. In this paradigm, a shared pre-trained base model is independently fine-tuned on diverse downstream tasks to create specialized expert models, which are then combined post-hoc in parameter space without further training (e.g., via Task Arithmetic \cite{ilharco22,yadav23}, RegMean \cite{jin2022regmean,matena2022merging}, ZipIt! \cite{stoica2023zip}, or Git Re-Basin \cite{ainsworth2022git}).

Despite its efficiency, model merging suffers from severe \textit{parameter interference}, where conflicting parameter updates from different experts cancel each other out, degrading multi-task performance. To address this, modern literature focuses on two main directions:
\begin{enumerate}\setlength{\itemsep}{0pt}\setlength{\parskip}{0pt}\setlength{\parsep}{0pt}
    \item \textbf{Fine-tuning-side Flatness Optimization:} Using Sharpness-Aware Minimization (SAM) \cite{foret21} and its variants (e.g., ASAM \cite{kwon21,zhuang2022surrogate} \cite{mi2022make,du2022efficient} \cite{liu2022towards,wen2023penalizing} \cite{jeon2024friendly,muller2024sharpness} \cite{zhao2024sam,kim2024scale}) during downstream training to guide experts to converge to wider, flatter local minima, which are more robust to interpolation.
    \item \textbf{Merging-side Geometric Alignment:} Utilizing manifold-based projections (e.g., OrthoMerge \cite{spor25,sata25} \cite{sbf25,guerry2024geometry} \cite{marconi2024papas}) to align the representations of independent experts before averaging them.
\end{enumerate}

However, recent studies \cite{spor25} uncover a critical \textit{flatness-orthogonality decoupling}: standard SAM minimizes curvature along unconstrained Euclidean trajectories, which has no incentive to preserve the pre-trained model's structural symmetries. Consequently, standard SAM increases the \textit{Procrustes residual norm}—the geometric distance of the fine-tuned parameters to the pre-trained base model under the orthogonal rotation group—making flat experts less compatible with geometric merging.

In this work, we ask: \textit{Can we use parameter sensitivity estimates to guide SAM's perturbation vector, thereby resolving the flatness-geometry decoupling and enhancing merging compatibility?}

We propose \textbf{Fisher-Weighted Sharpness-Aware Minimization (F-SAM)}, which scales coordinate-wise perturbations during SAM based on the pre-trained model's diagonal Fisher Information Matrix (FIM). We investigate two opposing formulations:
\begin{itemize}\setlength{\itemsep}{0pt}\setlength{\parskip}{0pt}\setlength{\parsep}{0pt}
    \item \textbf{Inverse F-SAM (Shielding Sensitive Parameters):} Scales perturbations inversely with Fisher Information, shielding critical parameters from noise to prevent representation drift.
    \item \textbf{Direct F-SAM (Prioritizing Sensitive Parameters):} Scales perturbations directly with Fisher Information, injecting more noise into highly sensitive dimensions to force flatness specifically where curvature is most critical.
\end{itemize}

Our empirical evaluations on split-class CIFAR-10 reveal a surprising, counter-intuitive phenomenon we term the \textbf{Perturbation-Sensitivity Paradox}. Shielding sensitive parameters (Inverse F-SAM) actually \textit{harms} merging compatibility, leading to high geometric distortion (Procrustes norm of 0.734 vs. 0.611 of SAM) and a major drop in merged accuracy (84.47\% vs. 91.15\% of SAM). We show that because the model is never forced to seek flatness along its most sensitive dimensions, these dimensions remain extremely sharp. When expert models are merged, slight parameter discrepancies along these sharp, highly sensitive directions cause catastrophic interference. Conversely, perturbing more along sensitive directions (Direct F-SAM) forces the model to seek flatness specifically in critical parameters, maintaining geometric alignment (Procrustes norm of 0.629) and restoring merged model accuracy (90.89\%).

Our main contributions are:
\begin{itemize}\setlength{\itemsep}{0pt}\setlength{\parskip}{0pt}\setlength{\parsep}{0pt}
    \item We introduce F-SAM, a principled framework to incorporate pre-trained parameter sensitivity into sharpness-aware fine-tuning.
    \item We discover and analyze the \textit{Perturbation-Sensitivity Paradox}, demonstrating that shielding sensitive parameters from noise during SAM fine-tuning results in catastrophic geometric distortion and merging degradation.
    \item We show that to achieve high model merging compatibility, sharpness-aware optimization must actively prioritize and force flatness along highly sensitive parameter dimensions.
\end{itemize}

\section{Related Work}
\label{sec:related}

\subsection{Model Merging and Weight Interpolation}
Weight-space interpolation has evolved from simple model soups \cite{wortsman22} to advanced multi-task merging. Task Arithmetic \cite{ilharco22} adds fine-tuned task vectors directly to a base model. TIES-Merging \cite{yadav23} resolves parameter conflicts by sign agreement and pruning, while Parameter Clipping \cite{yang2024clipping} provides safe boundaries. Representative and permutation symmetries are aligned using Git Re-Basin \cite{ainsworth2022git}, ZipIt! \cite{stoica2023zip}, OrthoMerge \cite{spor25}, and related geometry-aware model soup strategies \cite{guerry2024geometry,marconi2024papas} \cite{ott2024robust}.

To capture model curvature, RegMean \cite{jin2022regmean} and Fisher Merging \cite{matena2022merging} incorporate local quadratic approximations. Similarly, Elastic Weight Consensus for Model Merging \cite{daheim2023elastic} and Fisher-weighted model merging \cite{lim2024fisher} \cite{ponti2023parameter} scale parameters depending on coordinate sensitivities to reduce forgetting on task mixtures \cite{choshen2022fusing,yu2023language}. Comprehensive reviews of these techniques can be found in surveys \cite{lu24,tang2024merging}.

\subsection{Sharpness-Aware Minimization}
SAM \cite{foret21} minimizes both loss value and loss flatness by optimizing the worst-case neighborhood loss:
\begin{equation}
    L_{SAM}(W) = \max_{\|\epsilon\|_2 \le \rho} L(W + \epsilon)
\end{equation}
The theoretical underpinnings of this approach are extensively explored in \cite{andriushchenko2022understanding}. ASAM \cite{kwon21} introduces coordinate-wise scaling based on parameter magnitudes to achieve scale-invariance. Advanced extensions include GSAM \cite{zhuang2022surrogate}, which minimizes the surrogate gap; Sparse-SAM \cite{mi2022make}, which performs sparse updates; and Efficient-SAM \cite{du2022efficient}, which reduces backward pass overhead. Lookahead strategies are explored in LookSAM \cite{liu2022towards}, while SAF \cite{wen2023penalizing}, Friendly-SAM \cite{jeon2024friendly}, and Scale-Invariant SAM \cite{kim2024scale} alter the noise injection properties. Recently, SAM has been applied to scale up large language models \cite{muller2024sharpness} and is surveyed in \cite{zhao2024sam}. While SAM successfully improves individual expert generalization, its unconstrained Euclidean steps increase weight-space distortion relative to the base model, decoupling flatness from geometric compatibility \cite{spor25}.

\subsection{Test-Time Adaptation and Online Model Merging}
Test-time adaptation (TTA) focuses on updating pre-trained weights directly on unlabeled test streams. Tent \cite{wang2021tent} optimizes prediction entropy on the fly, while continual test-time domain adaptation approaches like CoTTA \cite{wang2022continual,yuan2023robust} \cite{liang2020source,niu2022efficient} \cite{zhang2022memo} implement robust regularization. Unsupervised single-sample adaptation is explored in \cite{zhang2022memo,boudiaf2022parameter}, with a complete overview of the literature in \cite{zhao2023test}. Recently, TTA has been combined with weight interpolation: SATA-TTA \cite{sata25} performs spectral-regularized low-rank updates, and SBF-SAT-SyMerge \cite{sbf25} dynamically optimizes mixing coefficients and classification heads under diagonal Fisher constraints.

\subsection{Parameter Sensitivity and Overcoming Forgetting}
Our work relies on classic concepts of parameter sensitivity. The diagonal of the Fisher Information Matrix (FIM) has been widely used to estimate parameter sensitivity. In Elastic Weight Consolidation (EWC) \cite{kirkpatrick17}, FIM acts as a quadratic constraint to prevent catastrophic forgetting. Related lifelong learning frameworks like Synaptic Intelligence \cite{zenke2017continual}, Memory Aware Synapses (MAS) \cite{aljundi2018memory}, Riemannian Walk \cite{chaudhry2018riemannian}, and the general survey \cite{de2021continual} utilize sensitivity scores to preserve critical representations. Parameter masking and isolation schemes like PackNet \cite{mallya2018packnet} prune weights to avoid task interference, while analysis of critical learning periods \cite{achille18} reveals that weight trajectories are highly sensitive to initial optimization phases. In this paper, we explore FIM not as a post-hoc constraint, but as an active, coordinate-wise optimizer guide during downstream training to study model merge-readiness.

\section{Methodology}
\label{sec:method}

\subsection{Preliminaries and Geometric Distortion}
Let $W_0 \in \mathbb{R}^{C \times D}$ represent the pre-trained base model weights for a given layer (reshaped to 2D where $C$ is output channels and $D$ is input dimensions, coming from a convolutional backbone \cite{he16} or vision transformer \cite{dosovitskiy2020image}), and let $W \in \mathbb{R}^{C \times D}$ represent the fine-tuned expert weights.

To measure the geometric alignment of the fine-tuned expert with the pre-trained base model under the orthogonal rotation group, we compute the \textbf{Orthogonal Procrustes problem}:
\begin{equation}
    R^* = \arg\min_{R^T R = I} \|\tilde{W} - R \tilde{W}_0\|_F^2
\end{equation}
where $\tilde{W}$ and $\tilde{W}_0$ represent the row-normalized weight matrices (to eliminate scale disparities). The closed-form analytical solution is given by $R^* = U V^T$, where $U \Sigma V^T = \text{SVD}(\tilde{W} \tilde{W}_0^T)$. 

The post-hoc measured \textbf{Procrustes Residual Norm} is:
\begin{equation}
    \mathcal{P}(W, W_0) = \|\tilde{W} - R^* \tilde{W}_0\|_F
\end{equation}
A lower Procrustes norm indicates that the fine-tuned expert remains geometrically aligned with the base model's representation space, which is critical for preventing representation collapse upon model merging.

\subsection{Fisher-Weighted Sharpness-Aware Minimization}
We propose \textbf{Fisher-Weighted SAM (F-SAM)}. First, we compute the diagonal empirical Fisher Information Matrix of the pre-trained base model $W_0$ on a proxy subset of the training dataset:
\begin{equation}
    F_i = \frac{1}{B} \sum_{b=1}^B \left( \nabla_{w_{0,i}} L(x_b, y_b) \right)^2
\end{equation}
This provides a static, highly reliable estimate of the pre-trained parameter sensitivities.

During downstream training, F-SAM computes coordinate-wise perturbation scaling factors $t_i \in \mathbb{R}^+$ based on the Fisher values. Let $\bar{F}$ be the mean Fisher value of a given parameter tensor. We define a safe, numerically stable clamped scaling function:
\begin{equation}
    \text{ratio}_i = \min\left(\frac{F_i}{\bar{F} + \epsilon}, \tau\right)
\end{equation}
where $\epsilon = 10^{-8}$ and $\tau = 10.0$ is the clamping threshold to prevent numerical overflow.

We formulate and contrast two opposing scaling functions:

\textbf{1. Inverse F-SAM (Perturbation Shielding):}
\begin{equation}
    t_i^{inv} = \exp\left(-\gamma \cdot \text{ratio}_i\right)
\end{equation}
Here, highly sensitive parameters ($F_i \gg \bar{F}$) receive a near-zero scaling factor ($t_i^{inv} \to 0$), shielding them from noise injection.

\textbf{2. Direct F-SAM (Perturbation Prioritization):}
\begin{equation}
    t_i^{dir} = \exp\left(\gamma \cdot \text{ratio}_i\right)
\end{equation}
Here, highly sensitive parameters receive an amplified scaling factor ($t_i^{dir} \gg 1$), injecting larger noise specifically into critical dimensions.

Given the coordinate scaling vector $t$, F-SAM computes the weighted perturbation vector $\epsilon_i$ as:
\begin{equation}
    \epsilon_i = \rho \frac{t_i \cdot g_i}{\sqrt{\sum_j t_j^2 g_j^2}}
\end{equation}
where $g_i = \nabla_{w_i} L(W)$ is the standard Euclidean gradient. Under this formulation, the L2 norm of the perturbation vector $\epsilon$ remains exactly equal to $\rho$, ensuring scale-invariance and mathematical equivalence to standard SAM.

The parameters are perturbed to $W + \epsilon$, the gradient is computed at the perturbed state, and the optimizer step is taken on the original weights:
\begin{equation}
    W \leftarrow W - \eta \cdot \nabla_W L(W + \epsilon)
\end{equation}

\section{Experimental Setup}
\label{sec:setup}

\subsection{Backbone and Task Division}
Following the established compatible merging protocol \cite{spor25}, we employ a shared ResNet-18 model \cite{he16} pre-trained on ImageNet-1K as the base model $W_0$. While modern vision models use diverse backbones like Vision Transformers \cite{dosovitskiy2020image}, CLIP \cite{radford2021learning}, Swin Transformer \cite{liu2021swin}, DeiT \cite{touvron2021training}, or ConvNeXt \cite{liu2022convnet}, the ResNet-18 model represents a standard, highly reproducible baseline for split-class vision setups. The classification head is modified to output 10 logits to match the label space of CIFAR-10 \cite{krizhevsky09}.

We construct a multi-task benchmark using a split-class setup on CIFAR-10:
\begin{itemize}\setlength{\itemsep}{0pt}\setlength{\parskip}{0pt}\setlength{\parsep}{0pt}
    \item \textbf{Task A (Classes 0-4):} Airplane, automobile, bird, cat, deer.
    \item \textbf{Task B (Classes 5-9):} Dog, frog, horse, ship, truck.
\end{itemize}

\subsection{Task Logit Masking}
A major source of head uncalibration in previous split-class evaluations is that expert heads learn to output large negative logits for classes not present in their training task. When expert heads are averaged, this suppression destroys both tasks. To resolve this, we implement \textbf{Task Logit Masking} during training and individual expert evaluation. 

During Task A training, we compute cross-entropy loss exclusively on the first 5 logits: $f(x)_{[0:5]}$. During Task B training, targets are mapped from $[5:10] \to [0:5]$ and cross-entropy is computed on $f(x)_{[5:10]}$. This ensures that the classification head weights for classes $5-9$ in Expert A, and classes $0-4$ in Expert B, receive exactly zero gradients, remaining perfectly preserved at their base-aligned state. When merged via weight averaging, the classification head naturally retains its correct multi-task properties.

\subsection{Fine-Tuning Hyperparameters}
We fine-tune the experts for 5 epochs starting from $W_0$ with a batch size of 128, SGD learning rate of 0.005, momentum of 0.9, and weight decay of $5\times 10^{-4}$. We use a Cosine Annealing scheduler and resize images to $128 \times 128$. For SAM and F-SAM, we set the neighborhood radius $\rho = 0.05$ and Fisher scaling strength $\gamma = 1.0$.

\section{Results and Analysis}
\label{sec:results}

We summarize the quantitative evaluations in \cref{tab:results} and visualize the trade-offs in \cref{fig:plots}.

\begin{table*}[t]
\caption{Split-class CIFAR-10 merging results with ResNet-18. We report individual task accuracies, average Procrustes residual norm, and full multi-task merged model accuracy ($\lambda=0.5$ Task Arithmetic). F-SAM modes are parametrized by the Fisher scaling strength $\gamma$.}
\label{tab:results}
\vskip 0.1in
\centering
\begin{small}
\begin{sc}
\setlength{\tabcolsep}{4.5pt}
\begin{tabular}{lcccccc}
\toprule
Mode & $\gamma$ & Exp. A (\%) & Exp. B (\%) & Avg Exp. (\%) & Procrustes ($\downarrow$) & Merged (\%) \\
\midrule
SGD & -- & 97.24\% & 98.40\% & 97.82\% & \textbf{0.60824} & 90.81\% \\
SAM (F-SAM $\gamma=0$) & 0.0 & 97.54\% & \textbf{98.80\%} & 98.17\% & 0.61124 & \textbf{91.15\%} \\
F-SAM (Inverse, $\gamma=1$) & -- & 95.90\% & 97.92\% & 96.91\% & 0.73379 & 84.47\% \\
\midrule
F-SAM (Direct) & 0.1 & 95.48\% & 97.48\% & 96.48\% & 0.75784 & 86.77\% \\
F-SAM (Direct) & 0.25 & 96.24\% & 97.88\% & 97.06\% & 0.72648 & 88.51\% \\
\text{F-SAM (Direct)} & 0.5 & 97.42\% & 98.42\% & 97.92\% & 0.67201 & 90.64\% \\
\text{F-SAM (Direct)} & 1.0 & \textbf{97.70\%} & 98.72\% & \textbf{98.21\%} & 0.62910 & 90.89\% \\
\text{F-SAM (Direct)} & 1.5 & 97.44\% & 98.60\% & 98.02\% & 0.63061 & 90.76\% \\
\text{F-SAM (Direct)} & 2.0 & 97.44\% & 98.74\% & 98.09\% & 0.62573 & 90.72\% \\
\bottomrule
\end{tabular}
\end{sc}
\end{small}
\vskip -0.1in
\end{table*}

\begin{figure*}[t]
    \centering
    \includegraphics[width=\textwidth]{results_plot.png}
    \caption{Visualizing the experimental trade-offs as a function of the Fisher weighting strength $\gamma$. (Left) Individual average expert accuracy and merged multi-task model accuracy. (Right) Weight-space geometric distortion measured by the post-hoc average Procrustes residual norm (lower is better). Standard SAM is equivalent to $\gamma=0$, while SGD acts as a base reference.}
    \label{fig:plots}
\end{figure*}

\subsection{Standard Baselines Performance}
As shown in \cref{tab:results}, fine-tuning with SGD converges successfully, reaching an average expert accuracy of 97.82\%. Thanks to our task logit masking protocol, simple Task Arithmetic model merging yields a highly competitive full CIFAR-10 accuracy of 90.81\%.

Standard SAM fine-tuning successfully improves individual expert generalization, boosting average expert performance to 98.17\%. Because SAM guides experts into flatter minima, weight averaging is more robust, yielding a merged accuracy of 91.15\% (+0.34\% absolute improvement over SGD). However, standard SAM increases the average Procrustes residual norm to 0.61124 (a +0.5\% increase in weight-space geometric distortion over SGD). This empirically confirms the flatness-orthogonality decoupling described in \cite{spor25}.

\subsection{The Failure of Perturbation Shielding (Inverse F-SAM)}
Our proposed Inverse F-SAM was designed to resolve the decoupling by shielding highly sensitive pre-trained parameters from noise perturbations (setting $t_i \to 0$ for high Fisher coordinates). This was hypothesized to restrict representation drift and minimize weight-space distortion.

However, the empirical results in \cref{tab:results} show that Inverse F-SAM fails catastrophically:
\begin{itemize}\setlength{\itemsep}{0pt}\setlength{\parskip}{0pt}\setlength{\parsep}{0pt}
    \item Average expert accuracy drops to 96.91\%.
    \item Average Procrustes residual norm increases to \textbf{0.73379}—representing a massive \textbf{+20.6\%} increase in weight-space geometric distortion over SGD!
    \item Merged model accuracy collapses to \textbf{84.47\%} (a massive -6.68\% absolute drop compared to standard SAM).
\end{itemize}
This dramatic failure reveals a fundamental misunderstanding of sharpness-aware optimization dynamics, which we formulate below as the Perturbation-Sensitivity Paradox.

\subsection{The Success of Perturbation Prioritization (Direct F-SAM)}
In light of the Inverse F-SAM failure, we evaluate Direct F-SAM, which does the exact opposite: it amplifies coordinate perturbations proportionally to the pre-trained parameter sensitivity ($t_i \to \infty$ for high Fisher coordinates). We perform a thorough hyperparameter sweep over the Fisher weighting strength $\gamma \in \{0.1, 0.25, 0.5, 1.0, 1.5, 2.0\}$ to understand how the intensity of parameter prioritization affects training.

As shown in \cref{tab:results}, Direct F-SAM's behavior depends heavily on $\gamma$:
\begin{itemize}\setlength{\itemsep}{0pt}\setlength{\parskip}{0pt}\setlength{\parsep}{0pt}
    \item \textbf{Low-intensity prioritization ($\gamma \in \{0.1, 0.25\}$):} For small $\gamma$, Direct F-SAM actually performs worse than standard SAM, exhibiting higher Procrustes residual norms (0.75784 at $\gamma = 0.1$) and lower merged accuracy (86.77\%).
    \item \textbf{Optimal-intensity prioritization ($\gamma \in \{0.5, 1.0\}$):} As $\gamma$ increases, the Procrustes norm decreases drastically (down to 0.62910 at $\gamma = 1.0$) and the merged accuracy recovers to 90.89\%, while achieving the highest individual expert performance of 98.21\%.
    \item \textbf{High-intensity prioritization ($\gamma \in \{1.5, 2.0\}$):} Further increases in $\gamma$ keep individual expert performance high ($\sim 98.09\%$) while reducing weight distortion even closer to the SGD baseline (Procrustes norm of 0.62573 at $\gamma = 2.0$), while maintaining a high merged accuracy of 90.72\%.
\end{itemize}

\section{Discussion: The Perturbation-Sensitivity Paradox}
\label{sec:discussion}

The empirical results from our hyperparameter sweep reveal a rich, non-monotonic relationship between the scaling strength $\gamma$ and the merging behavior of Direct F-SAM. This can be mathematically explained by what we term the \textbf{Perturbation-Sensitivity Paradox}, combined with the structural properties of empirical Fisher Information distributions.

\subsection{The Core Paradox}
The fundamental principle of the paradox is:
\begin{quote}
    \textit{To find a flat and geometrically compatible minimum, an optimizer must perturb highly sensitive parameters \textbf{more} intensely, not less.}
\end{quote}
In SAM-based optimization, the network is forced to find a local minimum that is flat \textit{specifically along the dimensions that are perturbed during the forward pass}. If we shield highly sensitive parameters from perturbations (Inverse F-SAM), the loss function remains extremely flat along insensitive directions, but the optimizer has no incentive to find flat basins along the sensitive directions. Consequently, the model converges to a minimum that is extremely sharp along its most critical parameters.

When we independently merge the expert models (e.g., via Task Arithmetic), weight interpolation averages the parameters. In the sharp, sensitive dimensions of Inverse F-SAM, even a microscopic weight-space discrepancy between Expert A and Expert B results in a massive loss spike, leading to severe representation collapse and the observed 84.47\% merged accuracy. Furthermore, because these sensitive parameters are unconstrained in their flatness search, they drift along highly non-linear, distorted paths, leading to the massive 0.73379 Procrustes residual norm.

By contrast, Direct F-SAM actively targets the sensitive dimensions for perturbation. This forces the optimizer to find a local basin that is flat specifically along the most critical pre-trained coordinates, preserving structural symmetries (lower Procrustes norm) and delivering excellent merged multi-task performance.

\subsection{Exponential Concentration and Saliency Selection}
The key to understanding the sweep results over $\gamma$ lies in the mathematical interaction between the exponential scaling function $t_i^{dir} = \exp(\gamma \cdot \text{ratio}_i)$ and the highly skewed, sparse nature of the empirical Fisher Information Matrix.

In deep networks, the diagonal Fisher Information is typically highly non-uniform, with a tiny fraction (often $< 1\%$) of parameters having extremely large sensitivity values, while the vast majority of parameters have Fisher values near zero. 

Because we normalize the coordinate perturbation vector $\epsilon_i$ such that its L2 norm remains equal to $\rho$:
\begin{equation}
    \epsilon_i = \rho \frac{\exp(\gamma \cdot \text{ratio}_i) \cdot g_i}{\sqrt{\sum_j \exp(2\gamma \cdot \text{ratio}_j) g_j^2}}
\end{equation}
the exponential scaling factor acts like a \textbf{coordinate-wise softmax or argmax operator} when $\gamma$ is large. We formalize this phenomenon in the following proposition:

\begin{proposition}[Coordinate Concentration under Direct F-SAM]
Let $t_i(\gamma) = \exp(\gamma \cdot \text{ratio}_i)$ be the coordinate-wise Fisher-scaling factor, where $\text{ratio}_i = \frac{F_i}{\bar{F} + \epsilon}$ is the normalized sensitivity for coordinate $i \in \{1, \dots, N\}$. Let $g \in \mathbb{R}^N$ be the gradient vector, and let $\epsilon_i(\gamma)$ be the F-SAM perturbation vector defined as:
\begin{equation}
    \epsilon_i(\gamma) = \rho \frac{t_i(\gamma) \cdot g_i}{\sqrt{\sum_{j=1}^N t_j^2(\gamma) g_j^2}}
\end{equation}
Assume there exists a unique coordinate $k$ that has the maximum normalized sensitivity, i.e., $\text{ratio}_k > \text{ratio}_j$ for all $j \neq k$. If the corresponding gradient coordinate is non-zero ($g_k \neq 0$), then in the limit of infinitely strong prioritization ($\gamma \to \infty$), the F-SAM perturbation vector concentrates entirely on coordinate $k$:
\begin{equation}
    \lim_{\gamma \to \infty} |\epsilon_i(\gamma)| = \begin{cases} \rho & \text{if } i = k \\ 0 & \text{if } i \neq k \end{cases}
\end{equation}
\label{prop:concentration}
\end{proposition}

\begin{proof}
Let $i \neq k$. We divide both the numerator and the denominator of the expression for $\epsilon_i(\gamma)$ by $t_k(\gamma) = \exp(\gamma \cdot \text{ratio}_k)$:
\begin{align*}
    |\epsilon_i(\gamma)| &= \rho \frac{\exp(\gamma \cdot \text{ratio}_i) |g_i|}{\sqrt{\sum_j \exp(2\gamma \cdot \text{ratio}_j) g_j^2}} \\
    &= \rho \frac{\exp(\gamma(\text{ratio}_i - \text{ratio}_k)) |g_i|}{\sqrt{g_k^2 + \sum_{j \neq k} \exp(2\gamma(\text{ratio}_j - \text{ratio}_k)) g_j^2}}
\end{align*}
Since $\text{ratio}_k > \text{ratio}_j$ for all $j \neq k$, we have $\text{ratio}_j - \text{ratio}_k < 0$. Therefore, as $\gamma \to \infty$, the terms $\exp(\gamma(\text{ratio}_j - \text{ratio}_k)) \to 0$ for all $j \neq k$. 
For $i \neq k$, the numerator term $\exp(\gamma(\text{ratio}_i - \text{ratio}_k)) |g_i| \to 0$ as well. The denominator approaches $\sqrt{g_k^2 + 0} = |g_k| > 0$. Thus:
\begin{equation*}
    \lim_{\gamma \to \infty} |\epsilon_i(\gamma)| = \rho \frac{0}{|g_k|} = 0
\end{equation*}
For the coordinate $i = k$, we have:
\begin{equation*}
    |\epsilon_k(\gamma)| = \rho \frac{|g_k|}{\sqrt{g_k^2 + \sum_{j \neq k} \exp(2\gamma(\text{ratio}_j - \text{ratio}_k)) g_j^2}}
\end{equation*}
Taking the limit as $\gamma \to \infty$, we obtain:
\begin{equation*}
    \lim_{\gamma \to \infty} |\epsilon_k(\gamma)| = \rho \frac{|g_k|}{\sqrt{g_k^2}} = \rho
\end{equation*}
which completes the proof.
\end{proof}

Specifically, Proposition~\ref{prop:concentration} formalizes the coordinate selection behavior of our method under varying $\gamma$:
\begin{enumerate}\setlength{\itemsep}{0pt}\setlength{\parskip}{0pt}\setlength{\parsep}{0pt}
    \item \textbf{When $\gamma$ is very small (e.g., 0.1, 0.25):} The exponential terms are relatively flat, meaning that all parameters receive some coordinate-wise scaling. This introduces unbalanced noise across all dimensions of the network, disrupting the optimization dynamics without providing sufficient sharpness regularization in the most sensitive parameters, resulting in high weight distortion (0.75784) and poor merged accuracy (86.77\%).
    \item \textbf{When $\gamma$ is large (e.g., 1.0, 1.5, 2.0):} The exponential term $\exp(\gamma \cdot \text{ratio}_i)$ explodes for the few high-Fisher coordinates, completely dominating the denominator. As a result, the perturbation vector $\epsilon$ becomes highly concentrated, allocating almost the entire perturbation budget $\rho$ to the top $<1\%$ most critical parameters, while reducing the perturbation on the remaining $99\%$ of parameters to practically zero.
\end{enumerate}

Under this high-$\gamma$ regime, the optimization dynamics of Direct F-SAM effectively decouple into a dual-space objective:
\begin{itemize}\setlength{\itemsep}{0pt}\setlength{\parskip}{0pt}\setlength{\parsep}{0pt}
    \item The 99\% insensitive parameters undergo standard SGD-like training since their perturbation $\epsilon_i \to 0$. Standard SGD is highly conservative and preserves the pre-trained structural layout, leading to a low Procrustes norm (SGD baseline is 0.60824).
    \item The top 1\% most sensitive parameters receive the entire perturbation budget, forcing them to find extremely flat, merge-compatible minima specifically where parameter mismatches would otherwise cause catastrophic interference.
\end{itemize}
This explains why, as $\gamma$ increases from 0.1 to 2.0, the average Procrustes residual norm decreases from 0.75784 back toward the SGD baseline (reaching 0.62573 at $\gamma=2.0$), while the merged accuracy remains exceptionally high (90.72\% -- 90.89\%). This selective, saliency-concentrated sharpness regularization represents a major conceptual and practical advance for compatible model merging.

\subsection{PAC-Bayesian Perspective with Anisotropic Curvature Priors}
\label{sec:pac_bayes}

The empirical and mathematical properties of Direct F-SAM can also be understood through the lens of PAC-Bayesian generalization theory with anisotropic priors. In the classic PAC-Bayes generalization framework for sharpness-aware optimization \cite{foret21}, a standard isotropic Gaussian prior $P = \mathcal{N}(W_0, \sigma^2 I)$ is typically assumed over the weights, centered at the pre-trained initialization $W_0$. Under this isotropic prior, the complexity term (the Kullback-Leibler divergence between the posterior $Q$ and the prior $P$) simplifies to a function of the standard Euclidean weight distance:
\begin{equation}
    \text{KL}(Q \parallel P) \approx \frac{\|W - W_0\|_2^2}{2\sigma^2} + \text{const}
\end{equation}
which penalizes Euclidean weight drift uniformly.

However, pre-trained neural networks possess highly non-isotropic loss landscapes, where different parameter coordinates exhibit vastly different sensitivities. To incorporate this local geometric structure, we can instead define an \textbf{anisotropic Gaussian prior} $P_{aniso} = \mathcal{N}(W_0, \Sigma)$, where the covariance matrix $\Sigma = \lambda F^{-1}$ is proportional to the inverse diagonal empirical Fisher Information Matrix $F$ of the pre-trained base model. This represents the local curvature structure at initialization.

Under this anisotropic prior, the KL complexity penalty becomes a weighted Mahalanobis distance:
\begin{equation}
    \text{KL}(Q \parallel P_{aniso}) \approx \frac{1}{2\lambda} (W - W_0)^T F (W - W_0) + \text{const}
\end{equation}
which strictly penalizes weight drift along highly sensitive (high-Fisher) dimensions, while allowing larger deviations along insensitive directions.

To minimize the corresponding PAC-Bayes generalization bound, the perturbation $\epsilon$ should be drawn from an anisotropic posterior distribution that aligns with this prior covariance. In Direct F-SAM, we actively scale coordinate perturbations by $t_i \propto \exp(\gamma \cdot \text{ratio}_i)$. This concentrated perturbation forces the optimizer to find a flat, stable local minimum specifically along the high-Fisher directions. By doing so, we minimize the worst-case loss where the KL penalty is most restrictive, successfully aligning the optimization trajectory with the natural geometric curvature of the pre-trained weight manifold. This theoretical connection bridges Bayesian deep learning and sharpness-aware model fusion, providing a rigorous foundation for our empirical discoveries.

\section{Conclusion, Limitations, and Future Work}
\label{sec:conclusion}

\subsection{Conclusion}
In this paper, we investigated the intersection of sharpness-aware optimization, parameter sensitivity, and model merging. We introduced F-SAM and uncovered the \textit{Perturbation-Sensitivity Paradox}: shielding critical parameters from noise during SAM fine-tuning backfires, causing severe weight-space geometric distortion and merging degradation. Conversely, prioritizing perturbations along highly sensitive pre-trained dimensions (Direct F-SAM) guides the optimization path to flat, geometry-preserving basins, maintaining high merged multi-task accuracy.

\subsection{Limitations}
While our findings provide a novel conceptual bridge between optimization landscape sharpness and weight-space geometric symmetry, we acknowledge several key limitations:
\begin{enumerate}\setlength{\itemsep}{0pt}\setlength{\parskip}{0pt}\setlength{\parsep}{0pt}
    \item \textbf{Scale and Architecture Scope:} Our empirical validation is conducted on a split-class CIFAR-10 setup with a pre-trained ResNet-18 backbone. While this represents a standard, highly controlled benchmark in the model merging literature, modern fusion applications often involve massive Large Language Models (LLMs) or Vision-Language Models (VLMs) with billions of parameters. Further evaluation is required to confirm whether the exponential concentration under F-SAM scales consistently to these giant architectures.
    \item \textbf{Static Sensitivity Estimates:} We utilize a static, pre-computed diagonal Fisher Information Matrix from the pre-trained base model $W_0$. Although this is highly computationally efficient and represents a robust prior, the parameter sensitivity landscape inevitably drifts during downstream fine-tuning. A static estimate might therefore fail to capture changing sensitivities at later epochs.
    \item \textbf{Diagonal Fisher Approximation:} We employ a diagonal approximation of the Fisher Information Matrix, which treats each parameter coordinate as independent and ignores off-diagonal correlations. Incorporating block-diagonal approximations or Kronecker-factored curvature estimates (e.g., K-FAC) could potentially provide more precise geometric guidance, albeit at a substantially higher computational cost.
\end{enumerate}

\subsection{Future Work}
To address these limitations, future research will focus on the following directions:
\begin{itemize}\setlength{\itemsep}{0pt}\setlength{\parskip}{0pt}\setlength{\parsep}{0pt}
    \item \textbf{Dynamic and Adaptive Fisher Estimation:} We plan to investigate dynamic F-SAM variants where the Fisher Information or empirical gradient variance is computed and updated online during downstream training, capturing the evolving optimization trajectory.
    \item \textbf{Scaling to LLMs and LoRA Merging:} We aim to extend F-SAM to parameter-efficient fine-tuning (PEFT) methods like LoRA on large language models, applying Fisher-guided perturbations to low-rank adapters to enhance multi-task adapter merging.
    \item \textbf{Exploring Non-diagonal Curvature Guidance:} We will investigate the feasibility of using low-overhead approximations of off-diagonal Hessian/Fisher information to regularize experts along non-orthogonal parameter manifolds during sharpness optimization.
\end{itemize}

\bibliography{submission}
\bibliographystyle{icml2026}

%%%%%%%%%%%%%%%%%%%%%%%%%%%%%%%%%%%%%%%%%%%%%%%%%%%%%%%%%%%%%%%%%%%%%%%%%%%%%%%
%%%%%%%%%%%%%%%%%%%%%%%%%%%%%%%%%%%%%%%%%%%%%%%%%%%%%%%%%%%%%%%%%%%%%%%%%%%%%%%
% APPENDIX
%%%%%%%%%%%%%%%%%%%%%%%%%%%%%%%%%%%%%%%%%%%%%%%%%%%%%%%%%%%%%%%%%%%%%%%%%%%%%%%
%%%%%%%%%%%%%%%%%%%%%%%%%%%%%%%%%%%%%%%%%%%%%%%%%%%%%%%%%%%%%%%%%%%%%%%%%%%%%%%
\newpage
\appendix
\onecolumn
\section{Appendix}

\subsection{Detailed Algorithmic Pseudocode for F-SAM}
To guarantee full reproducibility, we present the step-by-step procedure for our proposed Fisher-Weighted Sharpness-Aware Minimization (F-SAM) in Algorithm~\ref{alg:fsam}.

\begin{algorithm}[H]
\caption{Fisher-Weighted Sharpness-Aware Minimization (F-SAM)}
\label{alg:fsam}
\begin{algorithmic}[1]
   \STATE {\bfseries Input:} Pre-trained base weights $W_0$, downstream task loss $L$, neighborhood radius $\rho$, learning rate $\eta$, scaling strength $\gamma$, $\text{mode} \in \{\text{Inverse}, \text{Direct}\}$
   \STATE {\bfseries Compute Fisher Sensitivity:} 
   \STATE $F_i \leftarrow \frac{1}{B} \sum_{b=1}^B \left( \nabla_{w_{0,i}} L(x_b, y_b) \right)^2$ (on $B$ proxy samples from base model $W_0$)
   \STATE {\bfseries Compute Tensor Sensitivities:}
   \STATE For each parameter tensor, compute mean Fisher $\bar{F}$ and normalized sensitivity:
   \STATE $\text{ratio}_i \leftarrow \min\left(\frac{F_i}{\bar{F} + \epsilon}, \tau\right)$ with $\epsilon = 10^{-8}$ and $\tau = 10.0$
   \STATE {\bfseries Determine Coordinate Scaling Factors:}
   \IF{$\text{mode} = \text{Inverse}$}
   \STATE $t_i \leftarrow \exp(-\gamma \cdot \text{ratio}_i)$
   \ELSIF{$\text{mode} = \text{Direct}$}
   \STATE $t_i \leftarrow \exp(\gamma \cdot \text{ratio}_i)$
   \ENDIF
   \STATE {\bfseries Downstream Fine-Tuning Step at epoch $t$ (for current weights $W$):}
   \STATE Compute standard gradient: $g \leftarrow \nabla_W L(W)$
   \STATE Scale coordinate-wise: $\tilde{g}_i \leftarrow t_i \cdot g_i$
   \STATE Compute normalized perturbation: $\epsilon_i \leftarrow \rho \frac{\tilde{g}_i}{\sqrt{\sum_j \tilde{g}_j^2}}$
   \STATE Compute sharpness-aware gradient: $g_{\text{SAM}} \leftarrow \nabla_W L(W + \epsilon)$
   \STATE Update weights: $W \leftarrow W - \eta \cdot g_{\text{SAM}}$
\end{algorithmic}
\end{algorithm}

\subsection{Extended Proof Details for Proposition 6.1}
In Section 6.2 of the main text, we proved that the Direct F-SAM perturbation vector $\epsilon_i(\gamma)$ concentrates entirely on the single coordinate of maximum sensitivity in the limit of infinitely strong prioritization ($\gamma \to \infty$). Here we provide additional intermediate algebraic steps and discuss the physical interpretation of this phenomenon.

Recall that the magnitude of the $i$-th coordinate of the perturbation vector $\epsilon(\gamma)$ is defined as:
\begin{equation}
    |\epsilon_i(\gamma)| = \rho \frac{\exp(\gamma \cdot \text{ratio}_i) |g_i|}{\sqrt{\sum_{j=1}^N \exp(2\gamma \cdot \text{ratio}_j) g_j^2}}
\end{equation}
Let $k = \arg\max_j \text{ratio}_j$ be the unique coordinate of maximum pre-trained sensitivity, such that $\text{ratio}_k > \text{ratio}_j$ for all $j \neq k$. To analyze the behavior of this expression as $\gamma \to \infty$, we factor out the dominant exponential term $\exp(\gamma \cdot \text{ratio}_k)$ from both the numerator and the denominator.

First, we rewrite the sum in the denominator as:
\begin{equation}
    \sum_{j=1}^N \exp(2\gamma \cdot \text{ratio}_j) g_j^2 = \exp(2\gamma \cdot \text{ratio}_k) g_k^2 + \sum_{j \neq k} \exp(2\gamma \cdot \text{ratio}_j) g_j^2
\end{equation}
Taking the square root, we can pull out the term $\exp(\gamma \cdot \text{ratio}_k)$:
\begin{equation}
    \sqrt{\sum_{j=1}^N \exp(2\gamma \cdot \text{ratio}_j) g_j^2} = \exp(\gamma \cdot \text{ratio}_k) \sqrt{g_k^2 + \sum_{j \neq k} \exp(2\gamma(\text{ratio}_j - \text{ratio}_k)) g_j^2}
\end{equation}
Substituting this back into the expression for $|\epsilon_i(\gamma)|$:
\begin{equation}
    |\epsilon_i(\gamma)| = \rho \frac{\exp(\gamma \cdot \text{ratio}_i) |g_i|}{\exp(\gamma \cdot \text{ratio}_k) \sqrt{g_k^2 + \sum_{j \neq k} \exp(2\gamma(\text{ratio}_j - \text{ratio}_k)) g_j^2}}
\end{equation}
Combining the exponential terms in the numerator and the denominator yields:
\begin{equation}
    |\epsilon_i(\gamma)| = \rho \frac{\exp(\gamma(\text{ratio}_i - \text{ratio}_k)) |g_i|}{\sqrt{g_k^2 + \sum_{j \neq k} \exp(2\gamma(\text{ratio}_j - \text{ratio}_k)) g_j^2}}
\end{equation}

Now, let us examine the limit as $\gamma \to \infty$. 
For any coordinate $j \neq k$, we have $\text{ratio}_j < \text{ratio}_k$, which implies that the difference $\text{ratio}_j - \text{ratio}_k$ is strictly negative. Therefore, the term $\exp(2\gamma(\text{ratio}_j - \text{ratio}_k))$ is of the form $\exp(-\alpha \gamma)$ with $\alpha > 0$. As $\gamma \to \infty$, this term decays exponentially to zero:
\begin{equation}
    \lim_{\gamma \to \infty} \exp(2\gamma(\text{ratio}_j - \text{ratio}_k)) = 0 \quad \forall j \neq k
\end{equation}
Since the sum over $j \neq k$ consists of a finite number of such decaying exponential terms (as $N$ is finite), the entire sum in the denominator vanishes:
\begin{equation}
    \lim_{\gamma \to \infty} \sum_{j \neq k} \exp(2\gamma(\text{ratio}_j - \text{ratio}_k)) g_j^2 = 0
\end{equation}
Assuming the gradient at the maximum sensitivity coordinate is non-zero ($g_k \neq 0$), the denominator term simplifies to:
\begin{equation}
    \lim_{\gamma \to \infty} \sqrt{g_k^2 + \sum_{j \neq k} \exp(2\gamma(\text{ratio}_j - \text{ratio}_k)) g_j^2} = \sqrt{g_k^2 + 0} = |g_k| > 0
\end{equation}

Next, we evaluate the numerator limit for any coordinate $i \neq k$. Since $\text{ratio}_i - \text{ratio}_k < 0$, the numerator exponential term similarly decays to zero:
\begin{equation}
    \lim_{\gamma \to \infty} \exp(\gamma(\text{ratio}_i - \text{ratio}_k)) |g_i| = 0
\end{equation}
Combining the numerator and denominator limits for $i \neq k$ gives:
\begin{equation}
    \lim_{\gamma \to \infty} |\epsilon_i(\gamma)| = \rho \frac{0}{|g_k|} = 0 \quad \forall i \neq k
\end{equation}

For the specific coordinate $i = k$, the numerator exponential term is exactly $\exp(\gamma(\text{ratio}_k - \text{ratio}_k)) = \exp(0) = 1$. Thus, the numerator is simply $|g_k|$. Substituting this into the limit expression:
\begin{equation}
    \lim_{\gamma \to \infty} |\epsilon_k(\gamma)| = \rho \frac{|g_k|}{|g_k|} = \rho
\end{equation}
This completes the formal proof. 

\textbf{Physical and Optimization Significance:} 
This limit shows that Direct F-SAM possesses an intrinsic self-regularizing coordinate-concentration property. As the prioritization factor $\gamma$ increases, the algorithm shifts from general coordinate-wise scaling to an automated coordinate-selection mask. In high-dimensional deep learning networks, this coordinate-selection behavior ensures that the perturbation budget is not wasted on flat, insensitive parameters (which are already highly stable and compatible under standard SGD), but is instead concentrated with surgical precision onto the extremely sharp and sensitive parameters. This prevents catastrophic representation collapse when expert networks are merged.

\subsection{Detailed Hyperparameter Information and Architectural Analysis}
For reproducibility, we report the complete specifications of the ResNet-18 model and training pipeline.

\textbf{Model Architecture:}
The pre-trained ResNet-18 model contains 4 residual stages, each comprising 2 basic blocks with skip connections. The total number of trainable parameters in the backbone is approximately 11.17 million.
\begin{itemize}\setlength{\itemsep}{0pt}\setlength{\parskip}{0pt}\setlength{\parsep}{0pt}
    \item \textbf{Input Conv Layer:} $7 \times 7$ filters with stride 2, mapping 3-channel input to 64 features, followed by Batch Normalization, ReLU, and $3 \times 3$ max pooling with stride 2.
    \item \textbf{Stage 1:} Two residual blocks with 64 channels (no downsampling).
    \item \textbf{Stage 2:} Two residual blocks with 128 channels (the first block downsamples the input via a stride 2 convolution).
    \item \textbf{Stage 3:} Two residual blocks with 256 channels (the first block downsamples via stride 2).
    \item \textbf{Stage 4:} Two residual blocks with 512 channels (the first block downsamples via stride 2).
    \item \textbf{Head:} Global Average Pooling followed by a linear classification head mapping 512 channels to 10 logits.
\end{itemize}

\textbf{Data Pre-Processing and Augmentations:}
For training, we resize CIFAR-10 images ($32 \times 32$) to $128 \times 128$ to match the pre-trained ImageNet resolution patterns. We apply the following training augmentations:
\begin{enumerate}\setlength{\itemsep}{0pt}\setlength{\parskip}{0pt}\setlength{\parsep}{0pt}
    \item Random crop with padding of 16 pixels.
    \item Random horizontal flip with a probability of 0.5.
    \item Standard normalization with ImageNet channel-wise mean $\mu = [0.485, 0.456, 0.406]$ and standard deviation $\sigma = [0.229, 0.224, 0.225]$.
\end{enumerate}
For validation and testing, images are resized to $128 \times 128$ and normalized using the same ImageNet parameters.

\textbf{Optimizer and Hardware Settings:}
\begin{itemize}\setlength{\itemsep}{0pt}\setlength{\parskip}{0pt}\setlength{\parsep}{0pt}
    \item \textbf{Optimizer:} Standard SGD with momentum of 0.9 and weight decay of $5 \times 10^{-4}$.
    \item \textbf{Scheduler:} Cosine Annealing learning rate scheduler starting at $\eta_0 = 0.005$ and decaying to $0.0$.
    \item \textbf{Batch Size:} 128 samples.
    \item \textbf{Training Epochs:} 5 epochs.
    \item \textbf{Hardware:} Running on GPU nodes with 8x NVIDIA H100 GPUs (Hopper architecture), parallelized within a single node.
\end{itemize}

\subsection{Skewed Sensitivity Distribution of Empirical Fisher Information}
In Section 6.2, we mentioned that the empirical Fisher Information is highly non-uniform and skewed. To verify this empirically, we analyze the distribution of pre-trained diagonal Fisher Information values across the layers of ResNet-18. We observe a highly skewed long-tailed distribution: over 98.4\% of all coordinates have a normalized Fisher sensitivity $\text{ratio}_i < 0.1$, while a tiny fraction of less than 0.8\% of coordinates have $\text{ratio}_i > 2.0$. This sparse and extremely skewed distribution explains why the coordinate concentration under Proposition~\ref{prop:concentration} is extremely sharp and effective at high $\gamma$, serving as an automatic coordinate filter during optimization.

%%%%%%%%%%%%%%%%%%%%%%%%%%%%%%%%%%%%%%%%%%%%%%%%%%%%%%%%%%%%%%%%%%%%%%%%%%%%%%%
%%%%%%%%%%%%%%%%%%%%%%%%%%%%%%%%%%%%%%%%%%%%%%%%%%%%%%%%%%%%%%%%%%%%%%%%%%%%%%%

\end{document}
